\section{État de l'art}

\subsection{Concepts clés}
\paragraph{Forêt aléatoire.} Une forêt aléatoire (RF) agrège des arbres de décision entraînés sur des échantillons bootstrap, chaque scission ne considérant qu'un sous-ensemble aléatoire de variables \cite{breiman2001}. Ses hyperparamètres déterminants sont le nombre d'arbres, la profondeur maximale, la taille minimale des feuilles et le nombre de variables testées par scission.

\paragraph{Gestion du déséquilibre.} Deux familles s'opposent : le \emph{rééchantillonnage}, notamment SMOTE qui synthétise des exemples minoritaires par interpolation \cite{chawla2002}, et la \emph{pondération des classes}, qui pénalise davantage les erreurs sur les classes rares dans la fonction de coût ou le critère d'impureté (WRF \cite{chen2004,zhang2016}, régression logistique pondérée WLR \cite{maalouf2014}).

\paragraph{Optimisation d'hyperparamètres (HPO).} La recherche sur grille explore exhaustivement un produit cartésien de valeurs ; la recherche aléatoire tire les configurations dans des distributions et se montre plus efficace lorsque peu d'hyperparamètres sont influents \cite{bergstra2012} ; la recherche bayésienne (estimateur TPE) concentre les essais sur les régions prometteuses \cite{bergstra2011,akiba2019}.

\paragraph{Validation.} La validation croisée stratifiée à $K$ plis conserve les proportions de classes \cite{kohavi1995}. Lorsque la même validation sert à choisir les hyperparamètres \emph{et} à estimer la performance, l'estimation est optimiste ; la validation croisée imbriquée corrige ce biais \cite{cawley2010}. Enfin, la \emph{fuite de données} (\emph{leakage}) --- toute information de la cible ou de lignes de test disponible à l'entraînement --- gonfle artificiellement les scores \cite{kaufman2012}.

\subsection{Travaux antérieurs}
\textbf{Article source.} Bader \emph{et al.} \cite{bader2024} étudient 830 participants (dernière visite de chaque patient, parmi 2\,621 observations), décrits par 14 variables : cinq biomarqueurs de stress oxydatif (8-isoprostane, 8-OHdG, GSH, GSSG, rapport GSH/GSSG), l'haplotype ACE et huit variables cliniques ou démographiques. Six classifieurs (LR, RF, KNN, SVM, NB, ANN) sont comparés en classification binaire (PHQ-9\,$>4$) et à quatre classes (0--4, 5--13, 14--20, $\geq 21$), avec découpage 70/30 et validation croisée à 5 plis avec \emph{grid search}. La RF obtient une AUC de 0,84 (binaire) et 0,88 (quatre classes) avec toutes les variables ; la WRF atteint 0,91 à quatre classes, devant WLR et RF-SMOTE, la différence étant jugée significative par les tests de Friedman et de Conover. Le glutathion réduit (GSH) est la variable la plus importante par permutation.

\textbf{Approches concurrentes.} Sau et Bhakta \cite{sau2017} obtiennent une AUC de 0,94 avec une RF sur 520 personnes âgées à partir de variables socio-démographiques et de santé ; Priya \emph{et al.} \cite{priya2020} trouvent que la RF offre la meilleure exactitude sur des données déséquilibrées (échelle DASS-21) ; Zulfiker \emph{et al.} \cite{zulfiker2021} combinent SMOTE et AdaBoost (AUC de 0,96). Ces travaux rapportent surtout des performances, rarement l'espace de recherche des hyperparamètres ni la protection contre la fuite de données.

\subsection{Positionnement}
Notre travail \emph{reproduit le protocole} de \cite{bader2024} sur un jeu de données distinct (biomarqueurs mitochondriaux) puis \emph{l'améliore sur trois plans} : (a) une HPO explicite et comparée (grille, aléatoire, bayésienne), avec espaces publiés et coûts mesurés ; (b) une inférence statistique sur scores appariés issus d'une validation imbriquée, avec correction pour la dépendance entre plis \cite{nadeau2003} ; (c) un audit de fuite de données, absent des travaux précités, qui conditionne l'interprétation de tous les gains.
